בפרויקט זה לא נבנתה רשת עצבית אחת שמקבלת סרטון ומחזירה מיד את זהות כל
הרובוטים. במקום זאת, פותחה מערכת רב־שלבית המשלבת ממשק משתמש, שרת
עיבוד, שני מודלי זיהוי, אלגוריתם מעקב ומנגנון החלטה דטרמיניסטי.

\subsection{מבט־על על המערכת}

המערכת פועלת כצינור עיבוד סדרתי, שבו כל שלב מוסיף רובד מידע חדש ומצמצם
את אפשרויות הזיהוי, עד לקבלת החלטה דטרמיניסטית. בתרשים הבא ניתן לראות
את מבנה זרימת הנתונים בין הרכיבים השונים במערכת:

\begin{figure}[H]
  \centering
\begin{tikzpicture}[
  node distance=1.2cm and 1.5cm,
  base/.style={draw=black!70, thick, rectangle, rounded corners, align=center, font=\small, minimum height=0.8cm, inner sep=8pt},
  input/.style={base, fill=blue!10},
  model/.style={base, fill=green!10},
  logic/.style={base, fill=orange!10},
  data/.style={base, fill=yellow!15},
  output/.style={base, fill=purple!10},
  arrow/.style={-{Stealth[scale=1.2]}, thick, draw=black!80}
  ]

  \node (video) [input] {\texthebrew{קלט: סרטון וידאו} \\ \textenglish{(Full Match Feed)}};
  \node (yolo_robot) [model, below=of video] {\texthebrew{\textenglish{YOLO26}: זיהוי רובוטים} \\ \texthebrew{וסיווג בריתות}};
  \node (bytetrack) [model, below=of yolo_robot] {\texthebrew{ByteTrack: מעקב אובייקטים} \\ \textenglish{(Track IDs)}};

  \node (matchnum) [input, left=of video, xshift=-1.2cm] {\texthebrew{קלט: מספר מקצה}};
  \node (matchinfo) [data, below=of matchnum] {\texthebrew{שליפת קבוצות המקצה} \\ \textenglish{(Match Schedule Lookup)}};
  
  \node (yolo_digits) [model, right=of bytetrack, xshift=1.5cm] {\texthebrew{\textenglish{YOLO26}: זיהוי ספרות} \\ \texthebrew{(חיתוכי באמפרים)}};
  
  \node (voting) [logic, below=of bytetrack, yshift=-0.5cm] {\texthebrew{מנגנון הצבעה וזיהוי קבוצה} \\ \textenglish{(Levenshtein + Voting + Match Schedule)}};
  \node (smoothing) [logic, below=of voting] {\texthebrew{אינטרפולציה והחלקה גאוסית}};
  \node (final) [output, below=of smoothing] {\texthebrew{פלט: מסלולים מוחלקים} \\ \texthebrew{תצוגת \textenglish{Replay} באפליקציה}};

  \draw [arrow] (video) -- (yolo_robot);
  \draw [arrow] (yolo_robot) -- (bytetrack);

  \draw [arrow] (bytetrack.east) -- (yolo_digits.west)
    node[midway, above, font=\tiny] {\texthebrew{גזירת תמונות}};

  \draw [arrow] (yolo_digits.south) |- (voting.east)
    node[near start, right, font=\tiny, xshift=5pt] {\texthebrew{ספרות חזויות}};

  \draw [arrow] (bytetrack) -- (voting)
    node[midway, left, font=\tiny] {\texthebrew{נתוני מעקב}};

  % חצים של מספר המקצה ולוח המשחקים
  \draw [arrow] (matchnum) -- (matchinfo);
  \draw [arrow] (matchinfo.south) |- (voting.west)
    node[pos=0.25, left, font=\tiny] {\texthebrew{שש קבוצות אפשריות}};

  \draw [arrow] (voting) -- (smoothing);
  \draw [arrow] (smoothing) -- (final);

\end{tikzpicture}
  \caption{תרשים זרימת הנתונים במערכת.}
  \label{fig:system_pipeline}
\end{figure}

התרשים מתאר את זרימת הנתונים המרכזית במערכת בזמן הרצה. מבחינת המימוש
בפועל, ההרצה מתחילה באפליקציית \textenglish{Flutter}, שבה המשתמש בוחר
סרטון מקצה ומזין את מספר המקצה. האפליקצייה שולחת את הסרטון ואת מספר
המקצה לשרת \textenglish{FastAPI}, שמבצע את עיבוד הווידאו ואת כל שלבי
ה־\textenglish{inference}.

בשלב הראשון, השרת משתמש במספר המקצה כדי לשלוף את שש הקבוצות הרלוונטיות
מהלו"ז, שלוש לכל ברית. שימוש זה בידע מוקדם מצמצם את מרחב החיפוש מאלפי
קבוצות אפשריות לשש קבוצות בלבד, ומאפשר למערכת לשלב בין זיהוי חזותי
לבין מידע דומייני זמין מתחרות \textenglish{FRC}.

לאחר מכן הסרטון מעובד פריים אחר פריים. בכל פריים מופעל מודל
\textenglish{YOLO} לזיהוי רובוטים ולסיווג צבע הברית. תיבות הזיהוי
מועברות ל־\textenglish{ByteTrack}, שמחבר בין זיהויים בפריימים עוקבים
ויוצר \textenglish{Track ID} לכל רובוט לאורך זמן. עבור עקבות שעדיין
לא שויכו לקבוצה, השרת גוזר את אזור הרובוט ומפעיל עליו מודל
\textenglish{YOLO} נוסף לזיהוי הספרות שעל הבאמפר.

תוצאות זיהוי הספרות אינן משמשות כהחלטה מיידית על סמך פריים יחיד.
במקום זאת, הן מועברות לשכבת החלטה שמחברת בין צבע הברית, רשימת הקבוצות
במקצה, התאמת \textenglish{Levenshtein} והצבעה מצטברת לאורך זמן. שכבה
זו אחראית להפוך זיהויים חלקיים ורועשים להחלטה יציבה יותר לגבי זהות
הקבוצה של כל עקבה.

בסיום עיבוד הסרטון, המערכת מאגדת את נקודות המיקום של כל קבוצה, מבצעת
אינטרפולציה והחלקה, ומחזירה לאפליקצייה מבנה נתונים המתאר את מסלולי
הרובוטים לאורך המקצה. האפליקצייה משתמשת בפלט זה כדי להציג
\textenglish{Replay} אינטראקטיבי, כולל אפשרות לגלילה בציר הזמן ולהצגת
המסלולים מעל הסרטון המקורי.

ניתן לראות את המערכת כמחולקת לשלוש שכבות: שכבת הראייה הממוחשבת, שכבת
המעקב והשיוך, ושכבת הפלט ועיבוד המסלולים.

\subsection{מדריך למפתח: מבנה הקוד והרצת המערכת}

\subsubsection{מבנה כללי ומיפוי רכיבים לקוד}

מדריך המפתח מתבסס על תרשים זרימת הנתונים שהוצג באיור
\ref{fig:system_pipeline}. בעוד שהתרשים מציג את המערכת ברמת
הארכיטקטורה, חלק זה מתאר כיצד אותם רכיבים ממומשים בפועל בקבצי הקוד,
וכיצד ניתן להריץ את המערכת בסביבת פיתוח מקומית. בפרויקט הנוכחי, רוב
לוגיקת ה־\textenglish{backend} מרוכזת בקובץ \textenglish{main.py},
ולכן המיפוי מתייחס בעיקר לפונקציות המרכזיות בקובץ זה.

\begin{table}[H]
\centering
\begin{tabular}{|p{0.30\textwidth}|p{0.28\textwidth}|p{0.32\textwidth}|}
\hline
\textbf{רכיב} & \textbf{מימוש} & \textbf{תפקיד} \\ \hline

ממשק ושרת &
\textenglish{frontend/}, \textenglish{upload()}, \textenglish{status()} &
שליחת סרטון ומספר מקצה, יצירת משימת עיבוד וקבלת סטטוס. \\ \hline

Pipeline מרכזי &
\textenglish{process()}, \textenglish{run\_pipeline()} &
הרצת כל שלבי העיבוד על הסרטון והחזרת התוצאה הסופית. \\ \hline

זיהוי ומעקב &
\textenglish{model.track()}, \textenglish{frc\_bytetrack.yaml} &
זיהוי רובוטים וצבעי ברית, וחיבור זיהויים לאורך זמן בעזרת \textenglish{ByteTrack}. \\ \hline

OCR ושיוך קבוצות &
\textenglish{detect\_digits()}, \textenglish{update()} &
זיהוי ספרות מחיתוכי רובוטים, צבירת ראיות ושיוך עקבות לקבוצות. \\ \hline

פלט ו־Replay &
\textenglish{smooth\_and\_interpolate()} &
החלקת המסלולים והכנת מבנה נתונים להצגה באפליקצייה. \\ \hline

\end{tabular}
\caption{מיפוי בין רכיבי המערכת לבין המימוש בקוד.}
\end{table}

הטבלה אינה מפרטת כל פונקציה בקוד, אלא מציגה את הקשר בין שלבי המערכת
העיקריים לבין המימוש שלהם.

\subsubsection{מבנה תיקיות וקבצים מרכזיים}

הפרויקט מחולק לשתי תיקיות עיקריות: ה־\textenglish{backend} מרוכז
בתיקייה אחת, וה־\textenglish{frontend} הוא פרויקט
\textenglish{Flutter} נפרד. בתיקיית השורש נמצא גם
\textenglish{Makefile} שמאפשר להריץ את שני הרכיבים במקביל בזמן פיתוח.

\begin{english}
\begin{verbatim}
autoscout/
  Makefile

  backend/
    main.py
    requirements.txt
    start.sh
    frc_bytetrack.yaml

    models/
      robots-*.pt
      digits-*.pt

  frontend/
    ...
\end{verbatim}
\end{english}

שמות קובצי המודלים כוללים את סוג המשימה, גודל המודל ותאריך האימון.
לדוגמה, קובץ בשם \textenglish{robots-nano\_26\_05\_04.pt} מייצג מודל
לזיהוי רובוטים בגודל \textenglish{nano}, שאומן ב־4 במאי 2026. מבנה
שמות זה שימש במהלך הפיתוח כדי לעקוב אחרי גרסאות שונות של מודלים
וניסויים.

\begin{table}[H] \centering
  \begin{tabular}{|p{0.3\textwidth}|p{0.6\textwidth}|} \hline
    \textbf{רכיב} & \textbf{תפקיד} \\ \hline
    \textenglish{backend/main.py} & קובץ השרת המרכזי. כולל את טעינת
    המודלים, קבלת הווידאו, עיבוד הפריימים, הרצת \textenglish{YOLO},
    מעקב עם \textenglish{ByteTrack}, שיוך לקבוצות והחזרת תוצאות. \\
    \hline \textenglish{backend/models/} & תיקייה שבה נשמרים מודלי
    ה־\textenglish{YOLO} המאומנים: מודל זיהוי רובוטים ומודל זיהוי
    ספרות. \\ \hline
    \textenglish{backend/frc\_bytetrack.yaml} & קובץ הגדרות \textenglish{ByteTrack} שהותאם למעקב אחרי רובוטי \textenglish{FRC}. \\ \hline
    \textenglish{backend/requirements.txt} & רשימת ספריות \textenglish{Python} הנדרשות להרצת השרת. \\ \hline
    \textenglish{backend/start.sh} & סקריפט נוחות להפעלת השרת עם \textenglish{uvicorn}. \\ \hline
    \textenglish{frontend/} & פרויקט \textenglish{Flutter} שמציג את ממשק המשתמש: העלאת סרטון, מעקב אחרי עיבוד וצפייה ב־\textenglish{Replay}. \\ \hline
    \textenglish{Makefile} & מאפשר להריץ את ה־\textenglish{backend} ואת ה־\textenglish{frontend} במקביל מתוך תיקיית הפרויקט הראשית. \\ \hline
  \end{tabular}
  \caption{קבצים ותיקיות מרכזיים בפרויקט.}
\end{table}

\subsubsection{הרצת המערכת}

ניתן להריץ את רכיבי המערכת בנפרד, או להשתמש ב־\textenglish{Makefile}
להרצה משולבת בזמן פיתוח.

\noindent\textbf{הרצת ה־Backend}

בהרצה ראשונה יש ליצור סביבת \textenglish{Python} וירטואלית ולהתקין את
התלויות:

\begin{english}
\begin{verbatim}
cd backend
python -m venv .venv
source .venv/bin/activate
pip install -r requirements.txt
\end{verbatim}
\end{english}

לאחר האתחול הראשוני, ניתן להפעיל את השרת באמצעות סקריפט הנוחות:

\begin{english}
\begin{verbatim}
backend/start.sh
\end{verbatim}
\end{english}

הסקריפט עובר אוטומטית לתיקיית \textenglish{backend}, מפעיל את הסביבה
הווירטואלית \textenglish{.venv}, ומריץ את שרת ה־\textenglish{FastAPI}
באמצעות \textenglish{uvicorn} על פורט \textenglish{8000}.

\noindent\textbf{הרצה משולבת}

\begin{english}
\begin{verbatim}
make dev
\end{verbatim}
\end{english}

פקודה זו נועדה לנוחות בזמן פיתוח, והיא מפעילה את ה־\textenglish{backend}
ואת ה־\textenglish{frontend} במקביל. בפועל ניתן להריץ כל רכיב גם בנפרד.

\subsubsection{הערות מימוש ותחזוקה}

הקובץ \textenglish{frc\_bytetrack.yaml} מבוסס על הגדרות ברירת המחדל של
\textenglish{ByteTrack}, אך עבר התאמות למשימת המעקב אחרי רובוטים
בסרטוני \textenglish{FRC}. מטרת ההתאמות היא לשפר את יציבות
ה־\textenglish{Track IDs} במצבים של תנועה מהירה, חסימות חלקיות
ורובוטים שנמצאים קרוב זה לזה.

בגרסה הנוכחית של הפרויקט, רוב לוגיקת ה־\textenglish{backend} מרוכזת
בקובץ \textenglish{main.py}. מבנה זה נבחר משום שהמערכת ממוקדת יחסית:
קבלת סרטון, עיבודו, הרצת מודלי הזיהוי, מעקב, שיוך קבוצות והחזרת
פלט. ריכוז הקוד בקובץ אחד הקל על פיתוח מהיר, בדיקות ושינויים תכופים
במהלך העבודה.

בגרסה עתידית ניתן לשפר את תחזוקת הקוד על ידי פיצול
\textenglish{main.py} למודולים נפרדים, למשל:
\textenglish{inference.py} לטעינת המודלים והרצתם,
\textenglish{tracking.py} לניהול \textenglish{ByteTrack},
\textenglish{team\_assignment.py} לשכבת ההצבעה והשיוך,
ו־\textenglish{video\_processing.py} לעיבוד הווידאו.

\subsection{שכבת הראייה הממוחשבת}



מודל זיהוי הרובוטים מבוסס על ארכיטקטורת \textenglish{YOLO26n} ומודל
זיהוי הספרות על \textenglish{YOLO26s}. אלה מודלי ה־\textenglish{nano}
וה־\textenglish{small} --- הקטנים ביותר --- בסדרה החדשה ביותר של
\textenglish{YOLO}: בזכות כמות הפרמטרים הקטנה מתאפשרים ביצועים מהירים
גם על חומרה חלשה יותר, ובזכות מבנה ה־\textenglish{end-to-end} של סדרת
\textenglish{YOLO26}, אין צורך בשלב \textenglish{NMS} כחלק מהזיהוי, מה
שמשפר את המהירות אפילו יותר. בשביל מודל זיהוי הספרות נבחר מודל קצת
גדול יותר כי הוא מאפשר דיוק גבוה יותר בלי פגיעה משמעותית במהירות, כי
בין הספרות יש יותר מחלקות ויותר קושי להבדיל ביניהן.

בפרויקט זה נעשה שימוש בשני מודלים נפרדים המבוססים על ארכיטקטורת
\textenglish{YOLO26}. הטבלה הבאה מסכמת את ההבדלים המבניים והתפעוליים
ביניהם:

\begin{table}[H]
\centering
\caption{השוואת מפרט טכני: זיהוי רובוטים לעומת זיהוי ספרות.}
\label{tab:model_comparison}
\begin{tabular}{|l|c|c|}
\hline \textbf{מאפיין} & \textbf{מודל זיהוי הרובוטים} & \textbf{מודל
זיהוי הספרות} \\ \hline \textenglish{Base Model} &
\textenglish{YOLO26n} (Nano) & \textenglish{YOLO26s} (Small) \\ \hline
גודל קלט & $3 \times 640 \times 640$ & $3 \times 128 \times 128$
\\ \hline פרמטרים & 2.4M & 9.5M \\ \hline
שכבות & כ־122 & כ־329 \\ \hline \textenglish{FLOPs} & 5.4B & 5.18B
\\ \hline
\end{tabular}
\end{table}

הבחירה ב-\textenglish{YOLO26n} עבור זיהוי הרובוטים נועדה לאפשר קצב
פריימים גבוה בזמן אמת על תמונת המקור המלאה, בעוד שהשימוש
ב-\textenglish{YOLO26s} על קלט קטן ($128 \times 128$) אפשר דיוק גבוה
יותר בזיהוי הספרות תוך ניצול העובדה שהקלט כבר חתוך וקטן משמעותית.

\subsection{שכבת המעקב והשיוך}

\subsubsection{מעקב באמצעות \textenglish{ByteTrack}}

לאחר שלב הזיהוי, המערכת מפעילה \textenglish{ByteTrack} כדי לקשר בין
התיבות בפריימים עוקבים. כך, במקום להתייחס לכל זיהוי כאל אירוע נפרד,
המערכת בונה עקבות רציפות בזמן. לכל עקבה נשמר \textenglish{Track ID},
ובהמשך גם היסטוריית מיקומים של מרכז התיבה.

הבחירה ב־\textenglish{ByteTrack} מתאימה במיוחד לפרויקט זה, משום שהיא
מספקת איזון טוב בין מהירות לדיוק, ויכולה לעבוד היטב גם כאשר חלק מן
זיהויי הרובוטים חלשים יחסית. זה חשוב בסרטוני תחרות, שבהם יש טשטוש תנועה,
חסימות חלקיות ושינויים חדים בכיוון.

\subsubsection{מנגנון ההצבעה המצטבר}

החלק הייחודי ביותר בארכיטקטורה הוא שכבת ההחלטה. לכל \textenglish{Track
ID} נשמר וקטור ניקוד באורך 6, אחד לכל קבוצה אפשרית במקצה. בכל פעם
שמתקבל זיהוי ספרות חדש, המערכת מחשבת ניקוד חדש משני מקורות מידע:

\begin{itemize}
\item \textbf{צבע הברית:} אם הרובוט זוהה כברובוט כחול, רק שלוש הקבוצות
  של הברית הכחולה מקבלות יתרון, ולהפך עבור ברית אדומה.
\item \textbf{דמיון למספרי הקבוצות:} הספרות שזוהו מושוות לכל אחת משש
  הקבוצות במקצה. מכיוון שמודל ה־OCR עלול לטעות בספרה אחת עקב רעש
  ויזואלי או הסתרה, לא ניתן להסתמך על השוואת מחרוזות בינארית (נכון/לא
  נכון). במקום זאת, השתמשתי באלגוריתם מרחק לוינשטיין
  (\textenglish{Levenshtein Distance}), המחשב את מספר הפעולות המינימלי
  הנדרש להמרת מחרוזת אחת לאחרת. המרחק מנורמל, כך שמתקבלת החלטה רציפה —
  לדוגמה, זיהוי "3318" במקום "3316" יקבל ניקוד גבוה בזכות דמיון
  חלקי. פתרון זה מעניק למערכת חוסן (\textenglish{Resilience}) מפני
  שגיאות זיהוי נקודתיות.
\end{itemize}

במימוש הנוכחי שני מקורות המידע משוקללים באופן שווה, ולאחר מכן מעודכנים
בממוצע מעריכי לאורך הזמן. המשמעות היא שהחלטה סופית אינה מבוססת על פריים
בודד, אלא על הצטברות עקבית של ראיות. כך, אם זיהוי ספרות אחד שגוי, הוא
לא מפיל מייד את השיוך של העקבה כולה.

בנוסף, המערכת אוכפת אילוץ חד־חד־ערכי: בכל רגע קבוצה אחת יכולה להיות
משויכת רק לעקבה פעילה אחת. לכן, כאשר נשקל שיוך חדש, נלקחות בחשבון רק
קבוצות פנויות או קבוצות שהעקבה שלהן הוגדרה כ"אבודה" זמנית על ידי
העוקב.

\subsubsection{קבלת החלטה, סף ביטחון ומרווח דינמי}

שיוך סופי אינו מתקבל אוטומטית עבור המועמד בעל הניקוד הגבוה ביותר. כדי
לאשר זהות, המערכת דורשת שני תנאים: שהניקוד של המועמד המוביל יעבור סף
מינימלי, ושיהיה פער מספיק בינו לבין המועמד השני.

המרווח הזה אינו קבוע. אם שתי קבוצות במקצה בעלות מספרים דומים מאוד,
לדוגמה מספרים שנבדלים בספרה אחת בלבד, גם זיהוי ספרות טוב עשוי לתת
ניקוד דומה לשתי הקבוצות. לכן המרווח הנדרש בין המועמד המוביל לבין
המועמד השני מותאם לדמיון בין מספרי הקבוצות במקצה, שנמדד גם הוא באמצעות
\textenglish{Levenshtein} מנורמל.

כאשר מספרי הקבוצות שונים מאוד זה מזה, המערכת דורשת פער גדול יותר בין
המועמד הראשון לשני, משום שצפוי שיהיה קל יותר להבחין ביניהם. לעומת זאת,
כאשר מספרי הקבוצות דומים מאוד, המערכת מסתפקת בפער קטן יותר, משום שגם
זיהוי נכון עשוי להניב ציונים קרובים. עם זאת, המערכת עדיין דורשת
שהניקוד של המועמד המוביל יעבור סף מינימלי, ולכן לא מתקבלת החלטה על סמך
ראיה חלשה בלבד.

קטע הקוד לשלבי ההצבעה וקבלת ההחלטה צורף בקטע קוד
\ref{lst:voting_system} בנספח \ref{app:listings}.

\subsubsection{טיפול בחסימות ובהחלפת זהויות}

בסרטוני תחרות רובוטים נוטים לחצות זה את זה, להסתיר זה את זה ואף לצאת
לרגע מן הפריים. כדי לצמצם שגיאות זהות, המערכת מפעילה שני מנגנוני הגנה:

\begin{itemize}
\item אם עקבה שהייתה משויכת לקבוצה נעלמת ואינה נמצאת עוד בין העקבות
  הפעילות או האבודות, השיוך שלה משתחרר.
\item אם שתי עקבות שכבר קיבלו שיוך נראות כמי שחצו זו את זו בין שני
  פריימים עוקבים, המערכת מבטלת זמנית את האישור לשתיהן ומאפסת את ההצבעה,
  כדי למנוע מצב של החלפת זהויות שקטה.
\end{itemize}

מנגנונים אלו אינם מושלמים, אך הם משקפים בחירה ארכיטקטונית חשובה:
להעדיף מערכת שמסוגלת לחזור למצב של אי־ודאות ולבנות מחדש ביטחון, במקום
להיצמד לזהות שגויה לאורך זמן.

\subsection{שכבת הפלט ועיבוד המסלולים}

במהלך כל הריצה נשמר עבור כל \textenglish{Track ID} רצף של מיקומים,
המיוצג על ידי מרכז תיבת התיחום בכל חותמת זמן. לאחר מכן, ברגע שיש שיוך
לקבוצה, המערכת מאגדת יחד את כל קטעי העקבה המתאימים לאותה קבוצה, גם אם
במהלך המשחק הזהות בוטלה ואושרה מחדש.

כדי להפיק פלט שימושי לניתוח, מתבצעים שלושה שלבים נוספים:

\begin{itemize}
\item \textbf{איחוד ומיון כרונולוגי:} כל הנקודות של אותה קבוצה נאספות
  וממוינות לפי זמן.
\item \textbf{אינטרפולציה:} המסלול נדגם מחדש בקצב קבוע של
  \textenglish{\qty{50}{\Hz}}, כדי לקבל ייצוג אחיד ורציף.
\item \textbf{החלקה:} מופעל מסנן גאוסי חד־ממדי על צירי $x$ ו־$y$ כדי
  לצמצם רעש ולייצר מסלול חלק יותר לצפייה ולניתוח.
\end{itemize}

הפלט הסופי של השרת הוא מבנה נתונים הכולל את שמות הקבוצות, משך הסרטון,
גודל הפריים ורשימת נקודות עבור מסלול כל קבוצה. האפליקצייה משתמשת במידע
זה כדי להציג \textenglish{Replay}  אינטראקטיבי של המקצה, עם אפשרות להניח את המסלולים גם
מעל הווידאו המקורי.

\subsection{שיקולי תכנון והצדקה ארכיטקטונית}

הבחירה בארכיטקטורה זו נבעה מכמה שיקולים עיקריים:

\begin{itemize}
\item \textbf{ניצול ידע דומייני:} מספר המקצה מספק מידע חזק מאוד, ולכן
  חבל להתעלם ממנו ולפתור את הבעיה כאילו כל קבוצות \textenglish{FRC}
  הן מועמדות אפשריות בכל רגע.
\item \textbf{פירוק בעיה קשה לתת־בעיות פשוטות יותר:} זיהוי רובוטים,
  זיהוי ספרות, מעקב ושיוך הם אתגרים שונים. טיפול נפרד בכל אחד מהם מקל על
  האימון, על הניתוח ועל איתור שגיאות.
\item \textbf{עמידות לרעש:} זיהוי הספרות ברמת פריים בודד אינו מספיק
  אמין, ולכן נדרש מנגנון צבירת ראיות לאורך זמן במקום סיווג מיידי.
\item \textbf{יעילות חישובית:} מודל הספרות רץ רק על חיתוכים של עקבות
  שטרם קיבלו שיוך, ולא על כל רובוט בכל פריים, וכך נחסך זמן חישוב.
\item \textbf{הפרדה בין ממשק לעיבוד:} חלוקה ללקוח \textenglish{Flutter}
  ולשרת \textenglish{FastAPI} מאפשרת עיבוד אסינכרוני של סרטונים בלי
  לחסום את חוויית המשתמש.
\end{itemize}

לסיכום, ארכיטקטורת המערכת בפרויקט זה מבוססת על העיקרון של
\textbf{שילוב כמה רמזים חלשים להחלטה חזקה אחת}. במקום להסתמך על מודל
יחיד שיפתור את כל הבעיה בבת אחת, המערכת מחברת בין זיהוי, מעקב, ידע
מוקדם ולוגיקה של החלטה — וכך משיגה פתרון מעשי יותר לבעיית הסקאוטינג
האוטומטי.
